\section{Yoneda's Lemma}

\subsection{Representable functors}

\bdefn

    Let $\C$ be a category, then we have a bifunctor
    $$ \hom\colon\C^\op\times\C\to\Set $$
    call the {\emphcolor hom-functor} which maps pairs of objects $c,d$ to their hom-set $\hom(c,d)$.
    Pairs of morphisms $f\colon c\to d$ and $f'\colon c'\to d'$ are mapped to the function $\hom(f,f')\colon\hom(d,c')\to\hom(c,d')$ mapping $g\colon d\to c'$ to $f'\circ g\circ f\colon c\to d'$.
    I.e.\ $\hom(f,f')=f\bul f'$.

\edefn

The hom-functors restrict along $\C^\op$ and $\C$.
Specifically for any $c\in\C$ we have single-variable functors
$$ \hom(c,\bul)\colon\C\to\Set,\qquad \hom(\bul,c)\colon\C^\op\to\Set $$
where $\hom(c,f)=1_C\bul f$ which we will denote by $f_*$ and $\hom(c,f)=f\bul1_C=f^*$.

\bdefn

    A covariant functor $F\colon\C\to\Set$ is called {\emphcolor corepresentable} if there is some $X\in\C$ such that $F$ is naturally isomorphic to the hom-functor $\hom(X,\bul)$.
    A contravariant functor $F\colon\C^\op\to\Set$ is {\emphcolor representable} if $F$ is naturally isomorphic to $\hom(\bul,X)$ for some $X\in\C$.
    It is said that $F$ is {\emphcolor (co)represented} by $X\in\C$.

\edefn

We will sometimes say ``represented'' instead of ``corepresented'' when the variance is clear.

\bexam

    \benum
        \item The forgetful functor $U\colon\grp\to\Set$ is represented by $\bZ$.
        That is, we have a natural isomorphism $UG\cong\hom(\bZ,G)$ which maps $\phi\colon\bZ\to G$ to $\phi(1)$.
        This is a bijection (because $\bZ$ is the free group on a single generator), and it can be seen to be natural.

        \item In a somewhat generalization, the forgetful functor from the category of $R$-modules $U\colon\mod_R\to\Set$ is represented by $R$.
        Indeed, $\hom(R,M)\cong UM$ by mapping $\phi\mapsto\phi(1)$.

        \item The functor $O_n\colon\grp\to\Set$ mapping a group to the set of elements of order $\leq n$ is represented by $\bZ/n\bZ$.
        A morphism $\hom(\bZ/n\bZ,G)$ maps to its image on $1$ again.

        \item The contravariant powerset functor $\c P\colon\Set^\op\to\Set$, which maps sets to their powersets and $f\colon X\to Y$ to $f^{-1}\colon\c P(Y)\to\c P(X)$
        is represented by $\set{0,1}$.

        \item The open-set functor ${\rm Open}\colon{\sf Top}^\op\to\Set$ mapping a topology to its collection of open sets (and $f\colon X\to Y$ maps to $f^{-1}\colon{\rm Open}(Y)\to{\rm Open}(X)$)
        is represented by $\set{0,1}$ with the Sierpenski topology $\set{\emptyset,\Set0,\set{0,1}}$.
    \eenum

\eexam

Representations are useful due to the following lemma.

\subsection{Yoneda's lemma}

For two functors $F,G$ we denote by $\nat(F,G)$ the set of natural transformations $F\To G$ (i.e.\ $\nat$ is the hom-set of the functor category).

\bthrm[title=Yoneda's Lemma]

    Let $F\colon\C\to\Set$ be a functor and $X\in\C$.
    Then there is a bijection
    $$ \Phi\colon\nat(\hom(X,\bul),F) \xto{\ \sim\ } F(X) $$
    which is natural in both $X$ and $F$.

\ethrm

\bproof

    We build $\Phi$ explicitly: given a natural transformation $\alpha\colon\hom(X,\bul)\To F$, map it to $\alpha_X(1_X)$.
    To show its bijectivity, we explicitly construct its inverse as well.
    For $a\in F(X)$ let us define the natural transformation $\alpha\colon\hom(X,\bul)\To F$ whose components are
    $$ \alpha_Y\colon\hom(X,Y)\to F(Y),\qquad \alpha_Y(f) = F(f)(a) $$

    Explicitly, we have
    $$ \Phi(\alpha)=\alpha_X(1_X),\qquad \Psi(a)_Y(f) = F(f)(a) $$
    These are inverses:
    $$ \Phi(\Psi(a)) = \Psi(a)_X(1_X) = F(1_X)(a) = a,\qquad \Psi(\Phi(\alpha))_Y(f) = F(f)(\Phi(\alpha)) = F(f)\alpha_X(1_X) = \alpha_Y(f) $$
    where the final equality is due to naturality: $F(f)\alpha_X=\alpha_Yf_*$ so taking the value of this at $1_X$ gives $\alpha_Y(f)$.

    We show that $\Phi$ is natural in $X$ and $F$ separately; first we set $F$ constant and consider $\Phi_X$.
    To show that $\Phi_F$ is natural, let us consider what exactly our functors are.
    On the right clearly it is $F$, but on the left it is not as clear.
    Explicitly, we have a functor $\C\to[\C,\Set]^\op$ which maps $X\in\C$ to $h^X=\hom(X,\bul)$ and $f$ to $f_*\colon\hom(Y,\bul)\to\hom(X,\bul),h\mapsto hf$
    Since $\nat$ is the hom-functor $\nat(\bul,F)\colon[\C,\Set]^\op\to\Set$, we see that $\nat(h^\bul,F)$ is the composition:
    $$ \C \to [\C,\Set]^\op \xto{\nat(\bul,F)} \Set $$
    This maps $f\colon X\to Y$ to $f^{**}$ which maps $\alpha$ to $\alpha\circ f^*$ (whose $Z$ component maps $g\colon Y\to Z$ to $\alpha_Z(gf)$).

    So we want to show that $F(f)\circ\Phi_X=\Phi_Y\circ f^{**}$.
    For $\alpha\colon\hom(X,\bul)\To F$ the left maps to $F(f)(\alpha_X(1_X))$, and the right maps to $(\alpha f^*)_Y(1_Y)=\alpha_Y(f)$.
    We already showed (when showing $\Psi\Phi(\alpha)=\alpha$) that $F(f)\alpha_X(1_X)=\alpha_Y(f)$, as desired.

    Now we set $X$ and we consider $\Phi_F$ as $F$ varies.
    The left-hand side is simply the covariant hom-functor represened by $\hom(X,\bul)$, while the right-hand side is the functor ${\rm ev}_X\colon[\C,\Set]\to\Set$
    which evaluates $F\colon\C\to\Set$ at $F(X)$, and $\lambda\colon F\To G$ is mapped to ${\rm ev}_X(\lambda)=\lambda_X\colon F(X)\to G(X)$.
    Thus naturality just states that $\lambda_X\circ\Phi_F=\Phi_G\circ\lambda_*$ for all $\lambda\colon F\To G$.
    Indeed, if $\alpha\colon\hom(X,\bul)\To F$ is natural its image on the left is $\lambda_X\alpha_X(1_X)$ and its image on the right is $\Phi_G(\lambda\circ\alpha)=(\lambda_X\alpha_X)(1_X)$,
    which is equal.
    \qed

\eproof

Yoneda's lemma dualizes to:

\bcoro

    Let $F\colon\C^\op\to\Set$ be a contravariant functor and $X\in\C$, then there is a natural bijection
    $$ \nat(\hom(\bul,X),F) \cong FX,\qquad \alpha\mapsto\alpha_X(1_X) $$

\ecoro

In our proof we defined a functor $\C\to[\C,\Set]^\op$ which mapped $X\in\C$ to $\hom(X,\bul)$ and $f\colon X\to Y$ to $f_*\colon\hom(Y,\bul)\To\hom(X,\bul)$.
A functor $\C\to\D^\op$ is equivalent to a functor $\C^\op\to\D$ (with the same definition), so this is simply a functor $\C^\op\to[\C,\Set]$.
Replacing $\C$ with its opposite category, we get a functor $\C\to[\C^\op,\Set]$ defined as follows:
\benum
    \item $X\in\C$ is mapped to $\hom(\bul,X)$,
    \item $f\colon X\to Y$ is mapped to $f_*\colon\hom(\bul,X)\To\hom(\bul,Y)$ (composition with $f$).
\eenum

\bdefn

    The functor $\C\to[\C^\op,\Set]$ just described is called the {\emphcolor Yoneda embedding}, and it is denoted $\y$.
    We denote the related functor $\C^\op\to[\C,\Set]$ (the Yoneda embedding for the opposite category) by $\y^\op$ (or just $\y$ when no confusion arises).

\edefn

\bprop

    $\y$ is a fully faithful embedding.

\eprop

\bproof

    To show that $\y$ is fully faithful, we must consider for $X,Y\in\C$ the enduced function $\y_{XY}\colon\hom(X,Y)\to\nat(\y(X),\y(Y))$ and show it is a bijection.
    Notice that $\nat(\y(X),\y(Y))\cong\hom(X,Y)$ by the Yoneda lemma, which maps $\alpha\colon\y(X)\To\y(Y)$ to $\alpha_X(1_X)$.
    We know that this is a bijection, so it is sufficient to show that composition with $\y_{XY}$ is the identity (in either direction) in order to show that $\y_{XY}$ is a bijection.
    Let $f\colon X\to Y$, then $\y(f)_X(1_X)=f$ by definition, as desired.
    \qed

\eproof

Since fully faithful functors reflect isomorphisms:

\bprop

    For any two objects $X,Y\in\C$, $\hom(X,\bul)\cong\hom(Y,\bul)$ iff $X\cong Y$.
    Similarly for represented functors.
    Thus, a representation of a functor is unique up to isomorphism.

\eprop

\bdefn

    A contravariant functor $\C^\op\to\Set$ is called a {\emphcolor presheaf}.
    We denote the category of presheaves (i.e.\ $[\C^\op,\Set]$) by $\psh(\C)$.

\edefn

So the Yoneda embedding gives a fully faithful embedding $\C\to\psh(\C)$.

Recall that the center of a category is the commutative monoid of natural transformations $1_C\To1_C$.
Further recall that the center of equivalent categories are isomorphic.

\bexam

    We claim that $\grp$ and $\ab$ are not equivalent, we prove this by computing their centers.
    Recall that the forgetful functor $U\colon\grp\to\ab$ is represented by $\bZ$, and it is faithful.
    As it is faithful, it induces and embedding $Z(\grp)\to\nat(U,U)\cong\hom(\bZ,\bZ)\cong\bZ$.
    The bijection is given by $n\in\bZ$ maps to the natural transformations whose components are given by $x\mapsto x^n$.
    The only one of these natural transformations whose components are group morphisms are when $n=0,1$.
    Thus $Z(\grp)=\set{0,1}$ with multiplication.

    This shows that $Z(\ab)=\bZ$ since $x\mapsto x^n$ are all group morphisms.
    Thus $Z(\ab)\cong Z(\grp)$ and so they cannot be equivalent.

\eexam

\subsection{Representables and (co)limits}

\bprop

    Let $F\colon\C\to\Set$, then $F$ is corepresentable iff $\int F$ has an initial object.
    Dually $F\colon\C^\op\to\Set$ is representable iff $\int F$ has a terminal object.

\eprop

\bproof

    Suppose that $(c,x)\in\int F$ is initial.
    Recall the proof of the Yoneda lemma, in which we define the natural transformation $\Psi(x)\colon\hom(c,\bul)\To F(c)$ by $\Psi(x)_d(f)=Ff(x)$.
    For each $d\in\C$ and $y\in F(d)$ there is a unique $f\colon c\to d$ such that $Ff(x)=y$.
    Thus this $f$ is the unique $f\colon c\to d$ such that $\Psi(x)_d(f)=y$.
    Thus $\Psi(x)_d$ is a bijection for each $d\in\C$, and thus $\Psi(x)$ is a natural isomorphism represented $F$.

    Now suppose $\alpha\colon\hom(c,\bul)\To F$ is an isomorphism.
    Consider then $(c,\Phi(\alpha)=\alpha_c(1_c))\in\int F$: it is initial.
    Indeed, let $(d,x)\in\int F$; a morphism $(c,\Phi(\alpha))\to(d,x)$ is a morphism $f\colon c\to d$ such that $x=Ff(\Phi(\alpha))=\Psi(\Phi(\alpha))_d(f)=\alpha_d(f)$.
    Such a morphism exists and is unique since $\alpha_d\colon\hom(c,d)\to Fd$ is a bijection.
    Thus $(c,\Phi(\alpha))$ is initial.
    \qed

\eproof

Since limits are terminal objects in the category of elements of $\cone(\bul,F)$, we get the following:

\bcoro

    $F\colon I\to\C$ has a limit $\lim_IF$ if and only if $\lim_IF$ represents $\cone(\bul,F)$ (i.e.\ $\hom(\bul,\lim_IF)\cong\cone(\bul,F)$).
    Dually $F\colon I\to\C$ has a colimit $\colim_IF$ iff $\colim_IF$ corepresents $\cone(F,\bul)$.

\ecoro

Given a diagram $F\colon I\to\C$ and an element $X\in\C$, consider the functor $\hom(X,F\bul)\colon I\to\Set$.
Since $\Set$ is complete, this has a limit $\lim_I\hom(X,F\bul)$.
Recall that the limit of a diagram $F\colon I\to\Set$ is the equalizer of
$$ f,g\colon\prod_{i\in I}F(i)\tto\prod_{\alpha\colon i\to j}F(j),\qquad f=(p_j)_{\alpha\colon i\to j},\ g=(F(\alpha)p_i)_{\alpha\colon i\to j} $$
which is simply
$$ \lim_IG = \set{(a_i)_{i\in I}}[\forall\alpha\colon i\to j.\,a_j=F(\alpha)(a_i)] $$
So in our case
$$ \lim_I\hom(X,F\bul) = \set{(\lambda_i\colon X\to F(i))_{i\in I}}[\lambda_j=F(\alpha)\circ\lambda_i] $$
this is precisely the set of cones $X\To F$, i.e.\ $\lim_I\hom(X,F\bul)\cong\cone(X,F)$.
This isomorphism is natural.

Thus if $F$ has a limit $\lim_IF$ then it represents $\cone(\bul,F)$, i.e.\ 
$$ \lim_I\hom(X,F\bul) \cong \cone(X,F) \cong \hom(X,\lim_IF) $$
naturally.
This isomorphism also preserves the limit cone, so we have that $\hom(X,\bul)$ preserves limits.
Dually, $\hom(\bul,X)$ preserves limits in the opposite category, i.e.\ it carries colimits to limits.

\bprop

    The covariant hom-functors $\hom(X,\bul)$ preserve limits, while contravariant hom-functors $\hom(\bul,X)$ carry colimits to limits.

\eprop

Since isomorphic functors preserve the same limits, we get

\bprop

    Representable and corepresentable functors preserve limits (representable functors carry colimits to limits, as their codomain is the opposite category).

\eprop

